\documentclass[11pt]{article}

\usepackage[margin=1in]{geometry}
\usepackage{amsmath, amssymb, amsthm, bm, mathtools}
\usepackage{algorithm}
\usepackage{algpseudocode}
\usepackage{enumitem}
\usepackage{hyperref}

\hypersetup{
    colorlinks=true,
    linkcolor=blue,
    urlcolor=blue,
    citecolor=blue
}

\title{A Conditional Mutual Information Framework for\\
Interaction-Order Contributions in Learned Representations}
\author{J. P. Ibieta-Jimenez\\
\small Draft extension of the cross-fitted CMI framework}
\date{\today}

\newtheorem{definition}{Definition}
\newtheorem{theorem}{Theorem}
\newtheorem{proposition}{Proposition}
\newtheorem{lemma}{Lemma}
\newtheorem{corollary}{Corollary}
\newtheorem{remark}{Remark}

\newcommand{\E}{\mathbb{E}}
\newcommand{\I}{\mathbb{I}}
\newcommand{\R}{\mathbb{R}}
\newcommand{\1}{\mathbf{1}}

\begin{document}

\maketitle

\section{Introduction}

This document extends the cross-fitted conditional mutual information (CMI) framework from
\emph{scores} to \emph{interaction-order contributions} in learned representations.

The central question is:
\begin{quote}
How much unique predictive information about a target $Y$ is contributed by the order-$k$
interaction features, beyond what is already explained by the baseline representation and all
lower-order interactions?
\end{quote}

This allows us to move from score comparison to a principled information-theoretic analysis of
\emph{representation complexity}. In particular, it provides a formal way to evaluate whether
pairwise interactions are sufficient, whether higher-order interactions add predictive value, and
which interaction orders matter most.

\section{Setup and Notation}

Let
\begin{itemize}[leftmargin=1.5em]
    \item $Y \in \{0,1\}$ be a binary target,
    \item $R$ denote a baseline representation,
    \item $\Phi_1, \Phi_2, \ldots, \Phi_K$ denote feature blocks corresponding to interaction orders
    $1,2,\ldots,K$,
    \item $\Phi_{<k} := (\Phi_1,\ldots,\Phi_{k-1})$,
    \item $\Phi_{\le k} := (\Phi_1,\ldots,\Phi_k)$.
\end{itemize}

Interpretation:
\begin{itemize}[leftmargin=1.5em]
    \item $\Phi_1$ may contain unary / linear features,
    \item $\Phi_2$ may contain pairwise interactions,
    \item $\Phi_3$ may contain triplet interactions,
    \item and so on.
\end{itemize}

The role of the baseline representation $R$ is to prevent higher-order blocks from receiving credit
for predictive information that is already available in the underlying representation.

\subsection{The Role and Choice of the Baseline Representation $R$}

The choice of $R$ is one of the most consequential design decisions in the framework, because it
determines the precise meaning of the word ``unique'' in $\Delta_k$.
Concretely, $\Delta_k = I(Y;\Phi_k \mid R,\Phi_{<k})$ answers the question:
\emph{how much does $\Phi_k$ add, given that $R$ and all lower-order blocks are already known?}
Different choices of $R$ give genuinely different answers, each with a distinct scientific
interpretation.

\paragraph{$R = \varnothing$ (empty baseline).}
$\Delta_k$ measures the incremental value of $\Phi_k$ beyond the lower-order feature blocks alone.
This is the cleanest setting for synthetic validation, where no learned representation is
involved.

\paragraph{$R = $ a frozen backbone embedding.}
$\Delta_k$ measures how much $\Phi_k$ adds beyond what the backbone already captures.
This is the appropriate choice when evaluating whether explicit interaction features provide
complementary signal on top of a pretrained model.
Importantly, if the backbone was \emph{trained using $\Phi_k$} as part of its input or objective,
then $R$ and $\Phi_k$ are statistically entangled: conditioning on $R$ partially conditions on
$\Phi_k$ itself, leading to systematic underestimation of $\Delta_k$.
To avoid this circularity, $R$ should ideally come from a backbone trained
\emph{without explicit knowledge} of the $\Phi_k$ blocks being evaluated.

\paragraph{$R = $ the output of a specific intermediate layer $\ell$.}
This choice enables an \emph{interaction emergence profile}: by sweeping $\ell$ across depth, one
can track $\Delta_k(\ell) = I(Y;\Phi_k \mid \text{layer}_\ell,\Phi_{<k})$ and identify at which
depth the network ceases to benefit from the explicit interaction block.
At shallow layers, $\Delta_k(\ell)$ is expected to be large (the network has not yet encoded
order-$k$ structure); at deeper layers it should decay as the network learns to represent
interactions internally.

\begin{remark}[Effect of richer $R$ on estimator requirements]
As $R$ becomes higher-dimensional or more informative, the auxiliary model for
$P(Y \mid R, \Phi_{<k})$ must be flexible enough to use $R$ effectively.
A misspecified or underfit auxiliary model will underestimate $H(Y \mid R, \Phi_{<k})$,
inflating $\Delta_k$.
In practice, richer $R$ demands more expressive auxiliary estimators (e.g.\ gradient-boosted
trees rather than logistic regression), or careful dimensionality reduction of $R$ before
estimation.
\end{remark}

\section{Theoretical Foundation}

\subsection{Order-$k$ Conditional Mutual Information}

\begin{definition}[Order-$k$ Incremental Predictive Information]
The incremental predictive information carried by the order-$k$ interaction block is defined as
\begin{equation}
\Delta_k
:=
I(Y;\Phi_k \mid R,\Phi_{<k}).
\label{eq:delta_k_def}
\end{equation}
\end{definition}

This quantity measures the reduction in uncertainty about $Y$ when $\Phi_k$ is observed, given that
the baseline representation $R$ and all lower-order interaction blocks $\Phi_{<k}$ are already known.

\begin{remark}[Relativity of $\Delta_k$ to the choice of $R$]
$\Delta_k$ is not an absolute property of $\Phi_k$; it is a property of $\Phi_k$
\emph{relative to $R$}.
If $R$ already contains the information in $\Phi_k$ (e.g.\ because $R$ was produced by a model
trained with $\Phi_k$ as input), then $\Delta_k \approx 0$ not because $\Phi_k$ is uninformative
about $Y$, but because $R$ has already absorbed that information.
Conversely, if $R$ is an early-layer or task-agnostic embedding, $\Delta_k$ may be large even for
interaction orders that a deeper backbone would have implicitly learned.
The same $\Phi_k$ block can yield very different $\Delta_k$ values depending on which $R$ is used,
and comparing $\Delta_k$ across experiments is only meaningful when $R$ is held fixed.
\end{remark}

\subsection{Entropy Reduction Identity}

\begin{proposition}[Entropy Reduction Identity]
For each $k \ge 1$,
\begin{equation}
\Delta_k
=
H(Y \mid R,\Phi_{<k}) - H(Y \mid R,\Phi_{\le k}).
\label{eq:entropy_reduction_identity}
\end{equation}
\end{proposition}

\begin{proof}
By the definition of conditional mutual information,
\[
I(Y;\Phi_k \mid R,\Phi_{<k})
=
H(Y \mid R,\Phi_{<k}) - H(Y \mid R,\Phi_k,R,\Phi_{<k}).
\]
Since $(\Phi_k,\Phi_{<k}) = \Phi_{\le k}$, we obtain
\[
I(Y;\Phi_k \mid R,\Phi_{<k})
=
H(Y \mid R,\Phi_{<k}) - H(Y \mid R,\Phi_{\le k}).
\]
\end{proof}

\subsection{Interpretation}

Equation \eqref{eq:entropy_reduction_identity} shows that $\Delta_k$ is the
\emph{incremental conditional entropy reduction} obtained by adding the order-$k$ features to a
predictor that already has access to the baseline representation and all lower-order interactions.

Thus:
\begin{itemize}[leftmargin=1.5em]
    \item $\Delta_k > 0$ means order $k$ contributes unique predictive information,
    \item $\Delta_k \approx 0$ means order $k$ is conditionally redundant given lower orders,
    \item larger $\Delta_k$ means higher unique relevance of order $k$ for predicting $Y$.
\end{itemize}

\subsection{Binary Entropy Formulation}

For binary targets, define the binary entropy function
\begin{equation}
H_b(p) = -p \log p - (1-p)\log(1-p),
\qquad p \in (0,1).
\label{eq:binary_entropy}
\end{equation}

\begin{theorem}[Order-$k$ CMI via Binary Entropy Reduction]
For binary $Y$,
\begin{equation}
\Delta_k
=
\E\!\left[
H_b\!\left( P(Y=1 \mid R,\Phi_{<k}) \right)
-
H_b\!\left( P(Y=1 \mid R,\Phi_{\le k}) \right)
\right].
\label{eq:binary_entropy_reduction}
\end{equation}
\end{theorem}

\begin{proof}
Apply Proposition \eqref{eq:entropy_reduction_identity} and use the fact that for binary $Y$,
\[
H(Y \mid X) = \E\left[ H_b\!\left( P(Y=1 \mid X) \right) \right].
\]
Substituting $X=(R,\Phi_{<k})$ and $X=(R,\Phi_{\le k})$ yields the result.
\end{proof}

\section{Decomposition Across Interaction Orders}

\begin{theorem}[Chain Rule Decomposition Across Orders]
For any $K \ge 1$,
\begin{equation}
I(Y;\Phi_{\le K}\mid R)
=
\sum_{k=1}^K I(Y;\Phi_k \mid R,\Phi_{<k})
=
\sum_{k=1}^K \Delta_k.
\label{eq:chain_rule_orders}
\end{equation}
\end{theorem}

\begin{proof}
Apply the chain rule of conditional mutual information:
\[
I(Y;\Phi_1,\ldots,\Phi_K \mid R)
=
\sum_{k=1}^K I(Y;\Phi_k \mid R,\Phi_1,\ldots,\Phi_{k-1}).
\]
By definition, each summand is $\Delta_k$.
\end{proof}

Equation \eqref{eq:chain_rule_orders} motivates the \emph{interaction-order profile}
\[
(\Delta_1,\Delta_2,\ldots,\Delta_K),
\]
which quantifies how predictive information is distributed across interaction orders.

\section{Estimation via Cross-Entropy Reduction}

\subsection{Probabilistic Model Approach}

To estimate $\Delta_k$, we fit two probabilistic models:

\begin{align}
\hat p_{<k}(r,\phi_{<k})
&\approx
P(Y=1 \mid R=r,\Phi_{<k}=\phi_{<k}),
\label{eq:restricted_model}
\\
\hat p_{\le k}(r,\phi_{\le k})
&\approx
P(Y=1 \mid R=r,\Phi_{\le k}=\phi_{\le k}).
\label{eq:full_model}
\end{align}

The first is the \emph{restricted model}, which excludes order $k$.
The second is the \emph{full model}, which includes order $k$.

\subsection{Empirical Estimator}

Suppose we observe i.i.d. data
\[
\{(Y_i,R_i,\Phi_{1,i},\ldots,\Phi_{K,i})\}_{i=1}^n.
\]

\begin{definition}[Empirical Order-$k$ CMI Estimator]
The empirical estimator of $\Delta_k$ is
\begin{equation}
\hat \Delta_k
=
\frac{1}{n}
\sum_{i=1}^n
\left[
H_b\!\left( \hat p_{<k}(R_i,\Phi_{<k,i}) \right)
-
H_b\!\left( \hat p_{\le k}(R_i,\Phi_{\le k,i}) \right)
\right].
\label{eq:empirical_delta}
\end{equation}
\end{definition}

\subsection{Log-Loss Equivalence}

Define the binary cross-entropy loss
\begin{equation}
\ell(y,p) = -y\log p - (1-y)\log(1-p).
\label{eq:logloss}
\end{equation}

\begin{theorem}[Log-Loss Gain Equivalence]
The empirical estimator in \eqref{eq:empirical_delta} is equivalent to the average held-out
log-loss gain:
\begin{equation}
\hat \Delta_k
=
\frac{1}{n}
\sum_{i=1}^n
\left[
\ell\!\left(Y_i,\hat p_{<k}(R_i,\Phi_{<k,i})\right)
-
\ell\!\left(Y_i,\hat p_{\le k}(R_i,\Phi_{\le k,i})\right)
\right].
\label{eq:logloss_gain_equivalence}
\end{equation}
\end{theorem}

\begin{proof}
For binary targets, expected cross-entropy equals binary entropy when evaluated at the true
conditional probability. Thus, the entropy reduction formulation in
\eqref{eq:empirical_delta} is equivalent to reduction in expected log-loss, and empirically this is
estimated by the sample average in \eqref{eq:logloss_gain_equivalence}.
\end{proof}

Equation \eqref{eq:logloss_gain_equivalence} is the practical estimator we will use:
\[
\text{incremental CMI of order } k
\;\approx\;
\text{held-out log-loss of restricted model}
-
\text{held-out log-loss of full model}.
\]

\section{Cross-Fitting for Unbiased Estimation}

\subsection{Why Cross-Fitting is Necessary}

If the same data are used both to train $\hat p_{<k},\hat p_{\le k}$ and to evaluate
\eqref{eq:empirical_delta}, the resulting estimate is optimistically biased due to overfitting.
This is especially problematic in high-dimensional representation spaces.

\subsection{Cross-Fitted Estimator}

\begin{algorithm}[h]
\caption{Cross-Fitted Estimation of $\Delta_k = I(Y;\Phi_k\mid R,\Phi_{<k})$}
\begin{algorithmic}[1]
\Require Dataset $\mathcal{D}=\{(Y_i,R_i,\Phi_{1,i},\ldots,\Phi_{K,i})\}_{i=1}^n$, number of folds $M$, order $k$
\State Partition indices $\{1,\ldots,n\}$ into disjoint folds $F_1,\ldots,F_M$
\For{$m=1$ to $M$}
    \State $\mathcal{D}^{(m)}_{\mathrm{train}} \gets \{i : i \notin F_m\}$
    \State $\mathcal{D}^{(m)}_{\mathrm{test}} \gets \{i : i \in F_m\}$
    \State Train restricted model $\widehat P^{(m)}_{<k}$ on inputs $(R,\Phi_{<k})$
    \State Train full model $\widehat P^{(m)}_{\le k}$ on inputs $(R,\Phi_{\le k})$
    \For{$i \in F_m$}
        \State $\hat p^{(m)}_{<k,i} \gets \widehat P^{(m)}_{<k}(R_i,\Phi_{<k,i})$
        \State $\hat p^{(m)}_{\le k,i} \gets \widehat P^{(m)}_{\le k}(R_i,\Phi_{\le k,i})$
    \EndFor
    \State
    \[
    \widehat{\Delta}^{(m)}_k
    \gets
    \frac{1}{|F_m|}
    \sum_{i\in F_m}
    \left[
    \ell(Y_i,\hat p^{(m)}_{<k,i})
    -
    \ell(Y_i,\hat p^{(m)}_{\le k,i})
    \right]
    \]
\EndFor
\State
\[
\widehat{\Delta}^{\mathrm{CF}}_k
=
\frac{1}{M}\sum_{m=1}^M \widehat{\Delta}^{(m)}_k
\]
\State \Return $\widehat{\Delta}^{\mathrm{CF}}_k$
\end{algorithmic}
\end{algorithm}

\begin{remark}[Circularity risk when $R$ is a trained embedding]
The cross-fitting procedure eliminates overfitting bias in the auxiliary models
$\hat p_{<k}$ and $\hat p_{\le k}$.
However, it does \emph{not} resolve the deeper issue that arises when $R$ itself was trained on
data that includes $\Phi_k$: in that case, $R$ already encodes part of the signal in $\Phi_k$,
so the restricted model (which sees $R$) effectively has partial access to $\Phi_k$, and the
estimated $\Delta_k$ will be biased downward regardless of the cross-fitting design.
The remedy is to ensure that the backbone producing $R$ was trained on a held-out task or without
the explicit interaction features being evaluated.
In the synthetic validation experiments (Section~\ref{sec:synthetic}), this issue does not arise
because $R = \varnothing$.
\end{remark}

\subsection{Asymptotic Interpretation}

\begin{proposition}[Asymptotic Interpretation]
Under regularity conditions and sufficient calibration of the auxiliary predictors,
\begin{equation}
\widehat{\Delta}^{\mathrm{CF}}_k
\;\xrightarrow[n\to\infty]{}\;
I(Y;\Phi_k\mid R,\Phi_{<k}).
\label{eq:asymptotic_consistency}
\end{equation}
\end{proposition}

\begin{remark}
As in the original score-based framework, finite-sample bias depends on:
\begin{itemize}[leftmargin=1.5em]
    \item the quality of the auxiliary conditional probability estimators,
    \item model calibration,
    \item sample size relative to representation dimension,
    \item and the cross-fitting design.
\end{itemize}
Thus, $\widehat{\Delta}^{\mathrm{CF}}_k$ should be interpreted as a
\emph{predictive plug-in estimate of incremental conditional information}.
\end{remark}

\section{Order Selection and Adaptive Interaction Modeling}

\subsection{Order Relevance Profile}

The vector
\begin{equation}
\bm{\Delta}
=
(\Delta_1,\Delta_2,\ldots,\Delta_K)
\label{eq:delta_profile}
\end{equation}
defines the \emph{order relevance profile} of the task.

This profile answers:
\begin{itemize}[leftmargin=1.5em]
    \item Which interaction order first contributes substantial predictive information?
    \item Are pairwise interactions sufficient?
    \item Do triplet or higher-order terms add unique signal beyond lower orders?
\end{itemize}

\begin{remark}[Order profile is relative to $R$]
The order relevance profile $\bm{\Delta}$ characterizes task complexity
\emph{relative to the baseline $R$}.
A profile with large $\Delta_1$ and small $\Delta_k$ for $k \ge 2$ may reflect either that the
task is genuinely linear, or that $R$ has already captured all pairwise and higher-order structure.
When comparing profiles across architectures or tasks, $R$ must be held fixed to ensure
the comparisons are meaningful.
\end{remark}

\subsection{Learning Which Orders Matter}

Suppose we build a model with order-specific branches:
\[
\Psi(x) = \sum_{k=1}^K \alpha_k(x)\,\Phi_k(x),
\]
where $\alpha_k(x)$ are order gates.

The information-theoretic goal is to retain orders with high incremental value:
\[
\Delta_k = I(Y;\Phi_k \mid R,\Phi_{<k}).
\]

This suggests a principled order-selection rule:
\begin{itemize}[leftmargin=1.5em]
    \item keep order $k$ if $\Delta_k$ is large relative to its computational cost,
    \item prune order $k$ if $\Delta_k \approx 0$ given lower orders.
\end{itemize}

Thus, the CMI framework provides a formal target for \emph{adaptive interaction-order selection}.

\section{Specialization to Clifford-Style Order-2 Interactions}

For order $2$, let $\Phi_2$ denote the pairwise interaction block. In a Clifford-style architecture,
$\Phi_2$ may include both:
\begin{itemize}[leftmargin=1.5em]
    \item symmetric / coherence terms (inner-product-like),
    \item antisymmetric / structural terms (wedge-like).
\end{itemize}

Then
\[
\Delta_2 = I(Y;\Phi_2 \mid R,\Phi_1)
\]
measures the unique predictive information carried by pairwise relational structure beyond the
baseline representation and unary terms.

More finely, if $\Phi_2=(\Phi_2^{\mathrm{sym}},\Phi_2^{\mathrm{asym}})$, one may define
\begin{align}
\Delta_2^{\mathrm{sym}}
&=
I(Y;\Phi_2^{\mathrm{sym}} \mid R,\Phi_1),
\\
\Delta_2^{\mathrm{asym}}
&=
I(Y;\Phi_2^{\mathrm{asym}} \mid R,\Phi_1,\Phi_2^{\mathrm{sym}}),
\end{align}
to separate the predictive contribution of symmetric and antisymmetric pairwise structure.

\begin{remark}[Choice of $R$ for Clifford-style analysis]
In the order-2 case, the choice of $R$ has a direct impact on whether $\Delta_2$ captures
genuinely \emph{relational} structure or merely richer unary representations.
If $R$ is a strong unary encoder (e.g.\ the output of a deep MLP applied independently to each
input), then $\Delta_2 > 0$ provides clean evidence that pairwise structure adds information
beyond what per-element processing can capture.
If $R$ is weak or absent, $\Delta_2$ may absorb both unary and relational signal, making
interpretation less precise.
For the branch-wise decomposition into symmetric and antisymmetric terms, the ordering matters:
$\Delta_2^{\mathrm{sym}}$ is computed first, and $\Delta_2^{\mathrm{asym}}$ measures what remains
after the symmetric component is already known.
Swapping the order would yield different values and a different attribution.
\end{remark}

\section{Numerical Stability and Uncertainty Quantification}

\subsection{Clipped Binary Entropy}

For numerical stability, define the clipped binary entropy
\begin{equation}
H_b^{\varepsilon}(p)
=
H_b\!\left(\mathrm{clip}(p,\varepsilon,1-\varepsilon)\right),
\qquad \varepsilon > 0.
\label{eq:clipped_entropy}
\end{equation}

\subsection{Fold-Based Standard Error}

Given fold-specific estimates $\widehat{\Delta}_k^{(1)},\ldots,\widehat{\Delta}_k^{(M)}$, define
\begin{equation}
\mathrm{SE}\!\left(\widehat{\Delta}^{\mathrm{CF}}_k\right)
=
\sqrt{
\frac{1}{M-1}
\sum_{m=1}^M
\left(
\widehat{\Delta}^{(m)}_k
-
\overline{\Delta}_k
\right)^2
},
\qquad
\overline{\Delta}_k = \frac{1}{M}\sum_{m=1}^M \widehat{\Delta}^{(m)}_k.
\label{eq:fold_se}
\end{equation}

\section{Interpretation}

\begin{itemize}[leftmargin=1.5em]
    \item $\Delta_k > 0$: order $k$ provides unique predictive information beyond lower orders,
    \item $\Delta_k \approx 0$: order $k$ is conditionally redundant,
    \item larger $\Delta_k$: stronger evidence that the target depends on order-$k$ interactions.
\end{itemize}

Because the estimator is based on cross-entropy reduction, all quantities are measured in nats.
To convert to bits, divide by $\log 2$.

\section{Conclusion}

This framework generalizes cross-fitted conditional mutual information estimation from model
scores to interaction-order contributions in learned representations.

The core quantity,
\[
\Delta_k = I(Y;\Phi_k \mid R,\Phi_{<k}),
\]
measures the unique predictive information carried by order-$k$ interactions beyond the baseline
representation and all lower-order structure.

Its cross-fitted log-loss estimator provides a practical and principled way to:
\begin{itemize}[leftmargin=1.5em]
    \item quantify which interaction orders matter,
    \item compare pairwise vs higher-order interaction blocks,
    \item and guide adaptive interaction-order selection in neural architectures.
\end{itemize}

\section{Comparison Framework}

This section formalizes how the order-wise CMI quantities can be used to compare interaction
orders, architectural branches, and alternative representation families.

\subsection{Comparison of Interaction Orders}

For two interaction orders $j < k$, the quantities
\[
\Delta_j = I(Y;\Phi_j \mid R,\Phi_{<j}),
\qquad
\Delta_k = I(Y;\Phi_k \mid R,\Phi_{<k})
\]
measure different conditional contributions, since the conditioning sets differ.

Thus, $\Delta_j$ and $\Delta_k$ should not be interpreted as direct substitutes, but rather as
elements of an \emph{incremental hierarchy}:
\[
I(Y;\Phi_{\le K}\mid R) = \sum_{m=1}^K \Delta_m.
\]

The natural comparison questions are:
\begin{itemize}[leftmargin=1.5em]
    \item At which order does the first large gain in predictive information occur?
    \item Does the cumulative gain saturate at low order?
    \item Is the marginal contribution $\Delta_k$ negligible for sufficiently large $k$?
\end{itemize}

\begin{definition}[Cumulative Interaction Information]
For $k \ge 1$, define the cumulative interaction information up to order $k$ as
\begin{equation}
\Gamma_k
:=
I(Y;\Phi_{\le k}\mid R)
=
\sum_{m=1}^k \Delta_m.
\label{eq:cumulative_gamma}
\end{equation}
\end{definition}

The sequence $(\Gamma_1,\Gamma_2,\ldots,\Gamma_K)$ describes how predictive information accumulates
as interaction complexity grows.

\subsection{Comparison of Branches Within an Order}

Suppose order $k$ decomposes into branches,
\[
\Phi_k = (\Phi_k^{(1)},\ldots,\Phi_k^{(B)}).
\]

For example, in a Clifford-style order-2 block one may separate:
\[
\Phi_2 = (\Phi_2^{\mathrm{sym}}, \Phi_2^{\mathrm{asym}}),
\]
corresponding to symmetric/coherence and antisymmetric/structural terms.

We define branch-wise incremental contributions sequentially:
\begin{align}
\Delta_k^{(1)}
&:= I(Y;\Phi_k^{(1)} \mid R,\Phi_{<k}),
\\
\Delta_k^{(2)}
&:= I(Y;\Phi_k^{(2)} \mid R,\Phi_{<k},\Phi_k^{(1)}),
\\
&\vdots \nonumber \\
\Delta_k^{(B)}
&:= I(Y;\Phi_k^{(B)} \mid R,\Phi_{<k},\Phi_k^{(1)},\ldots,\Phi_k^{(B-1)}).
\label{eq:branchwise_cmi}
\end{align}

By the chain rule,
\begin{equation}
I(Y;\Phi_k \mid R,\Phi_{<k})
=
\sum_{b=1}^B \Delta_k^{(b)}.
\label{eq:branchwise_chain}
\end{equation}

This provides a principled decomposition of the total order-$k$ contribution into branch-specific
conditional gains.

\subsection{Comparison of Architectures}

Consider two representation families, $\mathcal{A}$ and $\mathcal{B}$, producing interaction blocks
\[
\Phi^{\mathcal{A}}_1,\ldots,\Phi^{\mathcal{A}}_K,
\qquad
\Phi^{\mathcal{B}}_1,\ldots,\Phi^{\mathcal{B}}_K.
\]

Examples include:
\begin{itemize}[leftmargin=1.5em]
    \item a Clifford-style order-2 architecture,
    \item a generic bilinear / quadratic interaction architecture,
    \item a tensorized higher-order interaction model,
    \item a standard FFN-based baseline.
\end{itemize}

For each architecture, define the order profile
\[
\bm{\Delta}^{\mathcal{A}} = (\Delta_1^{\mathcal{A}},\ldots,\Delta_K^{\mathcal{A}}),
\qquad
\bm{\Delta}^{\mathcal{B}} = (\Delta_1^{\mathcal{B}},\ldots,\Delta_K^{\mathcal{B}}).
\]

Then architecture comparison can be performed along at least three axes:
\begin{enumerate}[leftmargin=1.5em]
    \item \textbf{Total information captured:}
    \[
    \Gamma_K^{\mathcal{A}}
    \quad \text{vs.} \quad
    \Gamma_K^{\mathcal{B}}.
    \]
    \item \textbf{Order efficiency:} compare $\Delta_k$ per parameter or per FLOP.
    \item \textbf{Information concentration:} determine whether information is concentrated at low
    order or only appears at high order.
\end{enumerate}

\begin{remark}[Comparability requires a shared $R$]
Architecture comparison via order profiles is only meaningful when both architectures are
evaluated against the \emph{same} baseline $R$.
If $R^{\mathcal{A}} \neq R^{\mathcal{B}}$, differences in $\bm{\Delta}$ may reflect differences in
the baselines rather than in the interaction blocks themselves.
A natural choice is to use a shared, architecture-agnostic $R$ (e.g.\ raw input features, or a
frozen backbone trained independently of both architectures).
\end{remark}

\begin{definition}[Information Efficiency]
For an interaction block $\Phi_k$ with computational cost $C_k > 0$, define its information
efficiency as
\begin{equation}
\eta_k := \frac{\Delta_k}{C_k}.
\label{eq:information_efficiency}
\end{equation}
\end{definition}

This provides a cost-aware criterion for deciding whether an interaction order is worth retaining.

\section{Synthetic Validation Protocol}
\label{sec:synthetic}

Before applying the framework to learned representations on real datasets, the estimator should be
validated on synthetic tasks where the true relevant interaction order is known by construction.

\begin{remark}[Choice of $R$ in synthetic experiments]
Throughout this section, we set $R = \varnothing$ (the empty baseline).
This is the canonical choice for synthetic validation: there is no learned representation to
condition on, and the estimator's task is purely to identify the interaction order from the
raw feature blocks $\Phi_1, \Phi_2, \Phi_3$.
The empty baseline avoids the circularity and information-containment issues that arise with
learned $R$, and provides the clearest test of whether the CMI estimator correctly recovers
the ground-truth order structure.
\end{remark}

\subsection{Ground-Truth Task Families}

We recommend four classes of synthetic tasks:

\paragraph{Order-1 task.}
Let $X \in \R^d$ and define
\[
Y = \1\{w^\top X + \varepsilon > 0\}.
\]
The predictive structure is essentially linear / unary.

\paragraph{Order-2 task.}
Let
\[
Y = \1\{X_1 X_2 + \varepsilon > 0\}.
\]
The target depends on a pairwise interaction.

\paragraph{Order-3 task.}
Let
\[
Y = \1\{X_1 X_2 X_3 + \varepsilon > 0\}.
\]
The target depends on a triplet interaction beyond pairwise terms.

\paragraph{Mixed-order task.}
Let
\[
Y = \1\{a^\top X + b X_1 X_2 + c X_1 X_2 X_3 + \varepsilon > 0\},
\]
for nonzero coefficients $a,b,c$.

\subsection{Validation Goals}

The synthetic validation protocol should verify that:
\begin{itemize}[leftmargin=1.5em]
    \item for order-1 tasks, $\Delta_1$ dominates and $\Delta_k \approx 0$ for $k \ge 2$,
    \item for order-2 tasks, $\Delta_2$ is the first substantial gain,
    \item for order-3 tasks, $\Delta_3$ remains positive after conditioning on lower orders,
    \item for mixed-order tasks, the estimated order profile matches the construction.
\end{itemize}

\subsection{Estimator Diagnostics}

The following diagnostics should be reported:
\begin{itemize}[leftmargin=1.5em]
    \item mean and standard deviation of $\widehat{\Delta}_k^{\mathrm{CF}}$ across random seeds,
    \item fold-based uncertainty estimates,
    \item null distributions obtained via permutation of the candidate block $\Phi_k$,
    \item sensitivity to auxiliary model choice,
    \item sensitivity to sample size and dimensionality.
\end{itemize}

\section{Research Plan}

This section lays out the logical research program implied by the framework.

\subsection{Phase I: Stabilize the Estimation Backbone}

The first phase focuses on estimator reliability. The goal is to ensure that the cross-fitted
plug-in estimator
\[
\widehat{\Delta}^{\mathrm{CF}}_k
\]
is stable, interpretable, and capable of recovering the correct interaction order on controlled tasks.

\paragraph{Phase I objectives.}
\begin{itemize}[leftmargin=1.5em]
    \item Formalize the estimand $\Delta_k = I(Y;\Phi_k \mid R,\Phi_{<k})$ cleanly.
    \item Implement a robust cross-fitted estimator based on held-out log-loss reduction.
    \item Validate the estimator on synthetic tasks with known true order structure, using $R = \varnothing$.
    \item Quantify bias, variance, and ranking stability across seeds and folds.
\end{itemize}

\paragraph{Phase I deliverables.}
\begin{itemize}[leftmargin=1.5em]
    \item A short mathematical note (this document).
    \item A reusable estimator implementation.
    \item A synthetic benchmark suite.
    \item A calibration report showing when the estimator can and cannot be trusted.
\end{itemize}

\subsection{Phase II: Order-2 Representation Analysis}

The second phase uses the stabilized estimator to analyze structured order-2 interactions.

\paragraph{Phase II objectives.}
\begin{itemize}[leftmargin=1.5em]
    \item Compare unary vs pairwise representations.
    \item Decompose order-2 contributions into branch-wise components, e.g. symmetric vs
    antisymmetric.
    \item Determine whether order-2 structure contributes unique predictive information beyond the
    baseline representation.
\end{itemize}

\begin{remark}[Choice of $R$ in Phase II]
Phase II introduces a non-trivial $R$ for the first time.
The recommended approach is to use a \emph{frozen unary encoder} as $R$: a model trained using
only $\Phi_1$ (linear / unary features), so that $\Delta_2 = I(Y;\Phi_2 \mid R,\Phi_1)$
cleanly measures the additional value of pairwise structure beyond what a strong unary model
already captures.
If the analysis is layer-specific (e.g.\ evaluating at different depths of a neural network),
$R$ should be set to the output of the layer immediately preceding the interaction block, and
held fixed across all comparisons within Phase II.
The key invariant to maintain: $R$ must not have been trained using $\Phi_2$ as part of its input
or supervision signal.
\end{remark}

\paragraph{Phase II deliverables.}
\begin{itemize}[leftmargin=1.5em]
    \item An empirical comparison of order-1 and order-2 profiles.
    \item Branch-wise CMI decomposition.
    \item A case study on Clifford-style pairwise interactions.
\end{itemize}

\subsection{Phase III: Higher-Order Generalization}

The third phase extends the architecture and the analysis to higher-order interaction blocks.

\paragraph{Phase III objectives.}
\begin{itemize}[leftmargin=1.5em]
    \item Define practical feature constructions for $\Phi_3,\Phi_4,\ldots$.
    \item Estimate whether higher-order interactions add unique information beyond lower orders.
    \item Determine whether higher-order signal is statistically significant and computationally
    worthwhile.
\end{itemize}

\paragraph{Phase III deliverables.}
\begin{itemize}[leftmargin=1.5em]
    \item A hierarchy of higher-order feature blocks.
    \item Empirical estimates of $\Delta_3,\Delta_4,\ldots$.
    \item A comparison of cumulative information gain vs computational cost.
\end{itemize}

\subsection{Phase IV: Adaptive Order Selection}

The fourth phase uses the order-wise CMI quantities to guide architecture design.

\paragraph{Phase IV objectives.}
\begin{itemize}[leftmargin=1.5em]
    \item Build order-gated models with order-specific branches.
    \item Use $\widehat{\Delta}_k$ and $\eta_k$ to determine which orders should be retained.
    \item Study whether the interaction-order profile can be learned dynamically across tasks.
\end{itemize}

\paragraph{Phase IV deliverables.}
\begin{itemize}[leftmargin=1.5em]
    \item An adaptive-order interaction architecture.
    \item A cost-aware order-selection rule.
    \item A demonstration that the CMI framework can guide architectural pruning or expansion.
\end{itemize}

\section{Decision Roadmap: What Should Be Done Next?}

The framework naturally suggests several possible next steps. This section organizes them by logical
dependency.

\subsection{Option A: Improve the Estimator First}

This option prioritizes the statistical backbone.

\paragraph{Why choose this first?}
Because every later claim --- about order-2 usefulness, higher-order relevance, or adaptive-order
selection --- depends on the reliability of $\widehat{\Delta}_k$.

\paragraph{Recommended when:}
\begin{itemize}[leftmargin=1.5em]
    \item the main uncertainty is whether the current CMI estimation is stable enough,
    \item synthetic validation has not yet been completed,
    \item one wants a rigorous measurement tool before committing to architectural development.
\end{itemize}

\paragraph{Output:}
A trusted CMI estimation backbone.

\subsection{Option B: Analyze Order-2 Interactions First}

This option uses the current framework to study pairwise interactions in depth before moving to
higher order.

\paragraph{Why choose this first?}
Because order-2 is the simplest nontrivial case and already includes rich decompositions
(e.g. symmetric vs antisymmetric, inner vs wedge).

\paragraph{Recommended when:}
\begin{itemize}[leftmargin=1.5em]
    \item the estimator is already reasonably stable,
    \item one wants the quickest path to an interpretable empirical result,
    \item the goal is to connect the framework to Clifford-style interactions directly.
\end{itemize}

\paragraph{Output:}
A strong order-2 case study and possibly an early paper.

\subsection{Option C: Prototype Higher-Order Blocks Early}

This option prioritizes architectural exploration.

\paragraph{Why choose this first?}
Because higher-order interaction design may surface implementation bottlenecks, computational
constraints, or representational ideas that also inform the estimator.

\paragraph{Recommended when:}
\begin{itemize}[leftmargin=1.5em]
    \item the research goal is architecture-first,
    \item one expects representation design to be the main challenge,
    \item one is comfortable iterating with provisional estimators.
\end{itemize}

\paragraph{Output:}
A family of higher-order interaction models, though with weaker measurement guarantees early on.

\subsection{Recommended Logical Next Step}

The most logical next step is:

\begin{quote}
\textbf{Phase I: improve and validate the CMI estimator on synthetic tasks before moving to
higher-order architectural claims.}
\end{quote}

The reason is simple:
\begin{itemize}[leftmargin=1.5em]
    \item if the estimator is unstable, all downstream conclusions become fragile;
    \item if the estimator works well on known-order synthetic tasks, then later empirical claims
    about order-2, order-3, and adaptive-order selection become much more credible.
\end{itemize}

Thus the immediate next milestone should be:

\begin{center}
\framebox{
\parbox{0.9\textwidth}{
\textbf{Immediate milestone.}
Implement and benchmark the cross-fitted estimator
$\widehat{\Delta}^{\mathrm{CF}}_k$ on synthetic tasks with known ground-truth interaction order
(using $R = \varnothing$), and evaluate its stability across auxiliary model classes, folds,
seeds, and sample sizes.
}}
\end{center}

\section{Practical Milestone Checklist}

A concrete sequence for the next iteration is:

\begin{enumerate}[leftmargin=1.5em]
    \item Finalize this mathematical extension.
    \item Implement the order-$k$ estimator in reusable code.
    \item Build synthetic datasets with known order structure.
    \item Benchmark the estimator on these datasets (with $R = \varnothing$).
    \item Decide whether the estimator is reliable enough to support order-2 and higher-order
    empirical analysis.
    \item Only then, proceed to Clifford-style order-2 analysis and higher-order generalization,
    introducing non-trivial $R$ progressively and documenting its choice explicitly at each stage.
\end{enumerate}

\end{document}
